\chapter{Анализ задачи и обзор существующих решений}\label{ch:problem}

Цель главы "--- зафиксировать предметную область, в~которой строится
интеллектуальная система классификации текстовых документов, формально
описать решаемую задачу и~обосновать требования, на~которых строятся
архитектурные и~алгоритмические решения последующих глав. Глава
последовательно вводит контекст работы с~электронной почтой как
характерным классом текстовых документов, формализует задачу
многоклассовой классификации с~явной категорией \texttt{other},
рассматривает методы представления текста, классификации и~кластеризации,
сопоставляет известные системы организации входящей корреспонденции
и~переходит к~постановке задачи разработки.

\section{Предметная область и постановка проблемы}\label{sec:ch1/domain}

Электронная почта рассматривается в~работе как~представитель класса
текстовых документов с~большим потоком и~устойчивой потребностью
в~автоматическом упорядочении: объём ежедневной корреспонденции
профессионального пользователя измеряется десятками и~сотнями писем
\autocite{Singh2023EmFORE}, что делает ручную сортировку
по~папкам непродуктивной. Задача относится к~общей области
автоматической классификации текстов \autocite{Sebastiani2002TextCat},
однако обладает рядом особенностей, которые и~определяют постановку
настоящей работы.

Существующие инструменты эту постановку не~покрывают. Ручные правила
почтового клиента работают на~устойчивых признаках отправителя или
темы и~ломаются при~смене адресов и~форматов уведомлений. Встроенные
автоматические механизмы крупных почтовых сервисов "--- Focused~Inbox
в~Outlook \autocite{Microsoft2024FocusedInbox} и~Priority~Inbox
в~Gmail \autocite{Aberdeen2010PriorityInbox} "--- ограничены бинарной
шкалой «важное~/ неважное» и~не~позволяют пользователю задавать
собственные тематические или функциональные категории "--- проекты,
направления работы, типы документов "--- и~автоматически распределять
сообщения по~ним.

Из~этого следуют четыре содержательных свойства задачи, существенные
для дальнейших решений. Во-первых, набор классов \emph{персонализирован}
и~неизвестен заранее, поэтому система должна формировать его
по~исторической выборке писем, а~не~опираться на~готовый эталонный
корпус. Во-вторых, входной поток \emph{многоязычный (multilingual)}: русско-английская
переписка должна обрабатываться в~едином признаковом пространстве.
В-третьих, тексты писем содержат персональные данные и~фрагменты
внутренних документов, поэтому требуется \emph{локальная обработка}
без~передачи содержимого во~внешние сервисы. В-четвёртых, состав
классов и~распределение сообщений со~временем меняются, что относит
задачу к~постановкам с~явным дрейфом концепций.

Эти свойства согласуются с~двумя ключевыми перспективами в~прикладных
работах по~организации электронной почты. Подход TaskTracer показывает
применимость онлайн-классификаторов "--- наивного байесовского,
Passive-Aggressive и~Confidence-Weighted Linear Classification "---
для предсказания проектной принадлежности писем
\autocite{Keiser2009TaskTracer}. Подход EmFORE строит индуктивные
символьные правила распределения с~явной адаптацией к~изменению
пространства признаков и~дрейфу концепций \autocite{Singh2023EmFORE}.
Обе работы фиксируют общие требования: персонализированные классы,
обработка низкоуверенных предсказаний и~цикл обратной связи. На~этой
основе в~следующем разделе формулируется формальная постановка задачи.

\section{Формальная постановка задачи классификации}\label{sec:ch1/task}

\subsection{Многоклассовая постановка с категорией \texttt{other}}\label{subsec:ch1/task/formal}

Пусть \( X \) "--- множество допустимых текстовых сообщений, представленных
после нормализации в~виде канонической строки \texttt{normalized\_text}, и~пусть
\( Y_{\text{user}} = \{y_1, \dots, y_K\} \) "--- множество пользовательских
классов, формируемое в~процессе инициализации системы. Классификатор обучается
только этим классам; сообщения, которые не~удаётся уверенно отнести ни~к~одному
из~них, не~приписываются принудительно, а~получают \emph{отказ} от~классификации,
обозначаемый \texttt{other}. Решающая функция тем самым определяется
на~расширенном множестве исходов
\begin{equation}
    \label{eq:label-set}
    \widetilde{Y} = Y_{\text{user}} \cup \{\,\texttt{other}\,\},
\end{equation}
где \texttt{other} "--- не~самостоятельный класс, а~исход «воздержаться
от~решения».

Задача классификации формулируется как~построение отображения
\begin{equation}
    \label{eq:classifier}
    f \colon X \to \widetilde{Y},
\end{equation}
реализуемого как~композиция векторизатора \( \varphi \colon X \to \mathbb{R}^{d} \),
параметризованного классификатора \( f_\theta \colon \mathbb{R}^{d} \to \mathbb{R}^{K} \),
возвращающего апостериорные оценки принадлежности \( p_k(x) \) для каждого
пользовательского класса (классификатор \( f_\theta \) обучается только
классам \( Y_{\text{user}} \) и~метку \texttt{other} не~предсказывает),
и~решающего правила, учитывающего порог уверенности \( \tau \in (0, 1) \):
\begin{equation}
    \label{eq:decision-rule}
    f(x) =
    \begin{cases}
        \arg\max\limits_{k \in \{1, \dots, K\}} p_k(x), & \text{если } \max_{k} p_k(x) \geqslant \tau, \\
        \texttt{other},                                 & \text{иначе (отказ).}
    \end{cases}
\end{equation}
Параметры \( f_\theta \) настраиваются минимизацией эмпирического риска
на~размеченной выборке \( \{(x_i, y_i)\}_{i=1}^{N} \subseteq X \times Y_{\text{user}} \):
\begin{equation}
    \label{eq:empirical-risk}
    \widehat{R}(f_\theta) = \frac{1}{N} \sum_{i=1}^{N} \mathbb{1}\bigl[\,\hat{y}(x_i) \neq y_i\,\bigr],
    \qquad \hat{y}(x) = \arg\max_{k} p_k(x).
\end{equation}
Отказ \texttt{other} не~входит в~обучающую разметку и~не~участвует в~этой
величине: он определяется лишь порогом \( \tau \) на~этапе применения
правила~\eqref{eq:decision-rule}.

Отказ \texttt{other} выполняет две роли. Первая "--- безопасное поведение
при~низкой уверенности: при~\( \max_k p_k(x) < \tau \) система воздерживается
от~решения, что предпочтительнее ошибочного отнесения сообщения к~постороннему
классу. Вторая "--- исходы-отказы образуют отдельный поток сообщений, требующих
дальнейшего рассмотрения. В~общем случае правило~\eqref{eq:decision-rule}
допускает покомпонентные пороги \( \tau_k \) и~дополнительную проверку
принадлежности к~области известных данных (out-of-distribution) по~расстоянию
до~ближайшего центра кластера или по~калиброванной уверенности.

Общая структура постановки задачи показана на~рисунке~\ref{fig:classification-scheme}.

\begin{figure}[h]
    \centering
    \begin{tikzpicture}[
        thick,
        >=stealth,
        node distance=8mm and 10mm,
        block/.style={
                draw,
                rounded corners=2pt,
                align=center,
                minimum height=10mm,
                minimum width=28mm,
                inner sep=4pt,
                font=\small,
            },
        decision/.style={
                draw,
                diamond,
                aspect=2,
                align=center,
                inner sep=2pt,
                font=\small,
            },
        leaf/.style={
                draw,
                rounded corners=2pt,
                align=center,
                minimum height=9mm,
                minimum width=24mm,
                inner sep=4pt,
                font=\small,
                fill=black!4,
            },
    ]
    \node[block]                                    (msg)    {Сообщение \( x \in X \)\\(MIME, заголовки, тело)};
    \node[block, right=of msg]                      (norm)   {Нормализация\\\texttt{normalized\_text}};
    \node[block, right=of norm]                     (emb)    {Векторное\\представление \( \varphi(x) \)};
    \node[block, below=18mm of emb]                 (clf)    {Модель оценки\\уверенности\\\( g_\Theta \colon \mathbb{R}^{d} \to \mathbb{R}^{K} \)};
    \node[decision, left=of clf]                    (thr)    {\( \max_k s_k \geqslant \tau \)?};
    \node[leaf, left=of thr, yshift=10mm]           (label)  {Класс \( y_k \)\\\(+\)~оценка};
    \node[leaf, left=of thr, yshift=-10mm]          (other)  {\texttt{other}\\(буфер обратной связи)};

    \draw[->] (msg)  -- (norm);
    \draw[->] (norm) -- (emb);
    \draw[->] (emb)  -- (clf);
    \draw[->] (clf)  -- (thr);
    \draw[->] (thr.west) |- node[above, font=\scriptsize, pos=0.75] {да} (label.east);
    \draw[->] (thr.west) |- node[below, font=\scriptsize, pos=0.75] {нет} (other.east);
\end{tikzpicture}

    \caption{Общая постановка задачи классификации электронной почты
        с~явной категорией \texttt{other}: нормализация сообщения, построение
        векторного представления \( \varphi(x) \), вычисление оценок
        \( p_k(x) \) классификатором \( f_\theta \) и~применение порогового
        решающего правила}\label{fig:classification-scheme}
\end{figure}

\subsection{Особенности задачи}\label{subsec:ch1/task/specifics}

По~сравнению с~классической многоклассовой классификацией текстов решаемая задача
обладает рядом особенностей, связывающих свойства предметной области
(раздел~\ref{sec:ch1/domain}) с~введёнными объектами. Набор классов
\( Y_{\text{user}} \) не~задан заранее и~различается для разных пользователей, что
исключает эталонные корпуса и~требует автоматического формирования первичных классов
кластеризацией исторической почты с~ручной проверкой кандидатов (методы "--- далее
в~главе). Решающее правило~\eqref{eq:decision-rule} обязано допускать отказ
\texttt{other} для сообщений вне устоявшихся классов "--- разовой переписки, новых
уведомлений, смены структуры архива. Наконец, все компоненты "--- векторизатор
\( \varphi \) и~классификатор \( f_\theta \) "--- исполняются локально, без~передачи
текста писем во~внешние сервисы, а~\( \varphi \) отображает русско- и~англоязычные
сообщения в~единое признаковое пространство.

\subsection{Метрики качества}\label{subsec:ch1/task/metrics}

Для оценки качества классификатора используется набор стандартных
метрик многоклассовой классификации~\autocite{Sokolova2009Metrics,Manning2008IR},
дополненный показателями, характерными для постановки с~категорией
\texttt{other}.
Все метрики качества классификации вычисляются только по~пользовательским
классам \( Y_{\text{user}} \); отказ \texttt{other} не~является классом
и~отдельной метрикой качества в~работе не~оценивается.

Для каждого класса \( y_k \in Y_{\text{user}} \) на~тестовой выборке
определяются истинно положительные \( \mathrm{TP}_k \), ложно
положительные \( \mathrm{FP}_k \) и~ложно отрицательные \( \mathrm{FN}_k \)
предсказания. По~ним вычисляются покомпонентные точность и~полнота
\begin{align}
    \mathrm{Precision}_k & = \frac{\mathrm{TP}_k}{\mathrm{TP}_k + \mathrm{FP}_k}, \label{eq:precision} \\
    \mathrm{Recall}_k    & = \frac{\mathrm{TP}_k}{\mathrm{TP}_k + \mathrm{FN}_k}, \label{eq:recall}
\end{align}
их~гармоническое среднее
\begin{equation}
    \label{eq:f1}
    F_{1,k} = \frac{2\,\mathrm{Precision}_k \cdot \mathrm{Recall}_k}{\mathrm{Precision}_k + \mathrm{Recall}_k}
\end{equation}
и~агрегированные показатели:
\begin{equation}
    \label{eq:macro-f1}
    \mathrm{macro\text{-}F_1} = \frac{1}{K}\sum_{k=1}^{K} F_{1,k},
    \qquad
    \mathrm{Accuracy} = \frac{1}{N_{\mathrm{test}}} \sum_{i=1}^{N_{\mathrm{test}}} \mathbb{1}\bigl[\,\hat{y}_i = y_i\,\bigr].
\end{equation}

Сводное представление качества по~всем классам даёт матрица ошибок
(confusion matrix), которая используется как~для агрегированной оценки,
так и~для анализа систематических смешений между классами.

Помимо классических метрик фиксируются два операционных показателя, описывающих
объём отказов. Покрытие (coverage) "--- доля сообщений, для которых классификатор
принял решение в~пользу одного из~пользовательских классов:
\begin{equation}
    \label{eq:coverage}
    \mathrm{Coverage} = \frac{1}{N_{\mathrm{test}}}\sum_{i=1}^{N_{\mathrm{test}}} \mathbb{1}\bigl[\,f(x_i) \neq \texttt{other}\,\bigr].
\end{equation}
Доля отказов определяется как~\( 1 - \mathrm{Coverage} \). Оба показателя
характеризуют объём, а~не~качество отказа, и~рассматриваются совместно
с~macro-F1: при~фиксированном уровне ошибок более высокое покрытие означает
большую практическую применимость системы. Дополнительно анализируется
матрица ошибок, позволяющая выявлять систематические смешения пользовательских
классов.

\subsection{Дисбаланс классов и дрейф концепций}\label{subsec:ch1/task/imbalance}

Распределение пользовательских классов в~почтовом архиве заведомо
неравномерное: несколько активных проектов формируют большую часть
переписки, тогда как~ряд категорий представлен единичными примерами.
Из~показателей агрегированной оценки~\eqref{eq:macro-f1} accuracy в~таких
условиях маскирует деградацию качества на~редких классах, тогда как~macro-F1
усредняет качество по~классам равновесно; поэтому основной метрикой
сравнения моделей в~работе принят macro-F1.

Свойства потока писем меняются со~временем: появляются новые проекты
и~отправители, меняются форматы автоматических уведомлений, отдельные
классы теряют активность. Это явление известно как~дрейф концепций (concept drift)
\autocite{Gama2014ConceptDrift} и~характерно для рассматриваемой задачи;
в~смежной постановке оно представлено в~подходе EmFORE, реализующем явное
обновление правил классификации при~изменении распределения признаков
\autocite{Singh2023EmFORE}. Учёт дрейфа требует механизма обнаружения
расхождения активной модели с~текущим распределением писем и~последующего
обновления модели; соответствующие требования к~системе формулируются
в~разделе~\ref{sec:ch1/dev}.

\section{Методы представления текста для классификации}\label{sec:ch1/repr}

Векторизатор \( \varphi \colon X \to \mathbb{R}^{d} \) из~постановки~\eqref{eq:classifier}
переводит нормализованное сообщение в~векторное представление. На~этапе начального
формирования классов, когда размеченные примеры ещё отсутствуют, требуется единое
многоязычное представление, пригодное для~кластеризации исторической выборки и~её
ручной разметки; обученный затем на~полученных метках классификатор может опираться
на~собственное представление, настраиваемое непосредственно на~размеченных данных.
Выбор семейства \( \varphi \) определяет, какие из~свойств задачи, перечисленных
в~разделе~\ref{sec:ch1/domain}, могут быть выполнены принципиально, а~какие требуют
дополнительных усилий.

\subsection{Счётные представления}\label{subsec:ch1/repr/count}

Классическая отправная точка "--- модель «мешок слов» и~её взвешенный вариант
TF-IDF~\autocite{Manning2008IR,Salton1988TFIDF}. Терму~\( t \) в~сообщении~\( x \)
сопоставляется вес
\begin{equation}
    \label{eq:tfidf}
    w_{t,x} = \mathrm{tf}(t, x)\cdot \log\frac{N}{n_t},
\end{equation}
где \( w_{t,x} \) "--- вес терма~\( t \) в~сообщении~\( x \); \( \mathrm{tf}(t, x) \) "---
частота терма~\( t \) в~\( x \); \( N \) "--- число сообщений в~корпусе; \( n_t \) "---
число сообщений, содержащих~\( t \). Сообщение представляется разреженным вектором
\( \varphi_{\text{tfidf}}(x) \in \mathbb{R}^{|V|} \), где \( V \) "--- словарь корпуса.
Близость двух представлений измеряется косинусной мерой
\begin{equation}
    \label{eq:cosine}
    \cos\bigl(\varphi(x), \varphi(x')\bigr) =
    \frac{\langle \varphi(x), \varphi(x') \rangle}{\lVert \varphi(x) \rVert \, \lVert \varphi(x') \rVert},
\end{equation}
где \( \langle \cdot, \cdot \rangle \) "--- скалярное произведение, \( \lVert \cdot \rVert \) "---
евклидова норма; для единично нормированных векторов косинус совпадает со~скалярным
произведением и~монотонно связан с~евклидовым расстоянием.

Подход интерпретируем и~устойчив на~узких темах с~устоявшейся лексикой, но~не~учитывает
семантическую близость слов и~почти не~переносит признаки между языками: русские
и~английские формы одного понятия дают ортогональные измерения. Поэтому как~единственное
представление TF-IDF в~работе не~используется и~служит лишь для извлечения ключевых
терминов кластеров.

\subsection{Статичные дистрибутивные эмбеддинги}\label{subsec:ch1/repr/static}

Статичные дистрибутивные эмбеддинги "--- word2vec~\autocite{Mikolov2013word2vec},
GloVe~\autocite{Pennington2014GloVe}, fastText~\autocite{Bojanowski2017fastText}
(для русского "--- RusVectores~\autocite{Kutuzov2017RusVectores}) "--- фиксируют
дистрибутивную семантику отдельных слов, но~отображают слово в~единственный вектор
вне контекста. В~настоящей работе они не~используются и~приводятся как~историческая
ступень между счётными и~контекстными представлениями.

\subsection{Контекстные и многоязычные представления}\label{subsec:ch1/repr/context}

Современный слой методов "--- контекстные представления на~основе трансформеров.
BERT~\autocite{Devlin2019BERT} отдаёт зависимый от~контекста вектор каждой позиции,
а~кросс-языковое расширение XLM-R~\autocite{Conneau2020XLMR} обучается совместно
на~ста языках, включая русский. Для~кластеризации и~разметки удобнее однородные
представления уровня предложения (Sentence-BERT~\autocite{Reimers2019SBERT}), причём
многоязычный вариант получают дистилляцией знаний из~монолингвального
энкодера~\autocite{Reimers2020MultilingualSBERT} либо сквозным обучением на~параллельных
предложениях (LaBSE~\autocite{Feng2022LaBSE}). Контекстный энкодер отображает сообщение
в~единый вектор усреднением представлений токенов с~последующей единичной нормировкой:
\begin{equation}
    \label{eq:sentence-embedding}
    \varphi(x) = \frac{1}{\lvert T(x) \rvert} \sum_{t \in T(x)} \mathbf{e}_t,
    \qquad \widehat{\varphi}(x) = \frac{\varphi(x)}{\lVert \varphi(x) \rVert},
\end{equation}
где \( T(x) \) "--- множество (непустых) позиций токенов сообщения~\( x \);
\( \lvert T(x) \rvert \) "--- их число; \( \mathbf{e}_t \in \mathbb{R}^{d} \) "---
контекстный вектор токена~\( t \), выдаваемый трансформерным энкодером;
\( \widehat{\varphi}(x) \) "--- нормированное представление, для которого косинусная
близость~\eqref{eq:cosine} сводится к~скалярному произведению. Такое представление
уровня предложения покрывает требование многоязычности и~отдаёт векторы, пригодные
для нелинейной классификации и~плотностной кластеризации; именно оно используется
в~работе для начального формирования классов.

\subsection{Сводная типология}\label{subsec:ch1/repr/taxonomy}

Описанные семейства образуют последовательность с~растущей выразительной
способностью и~усложняющимися условиями обучения и~исполнения
(рисунок~\ref{fig:embedding-taxonomy}). Счётные и~статичные представления остаются
базовой линией, однако только контекстные многоязычные модели уровня предложения
одновременно отвечают требованию многоязычности и~не~требуют переобучения с~нуля
на~локальном архиве; поэтому они приняты как~основное представление для начального
формирования классов.

\begin{figure}[h]
    \centering
    \resizebox{\textwidth}{!}{\begin{tikzpicture}[
        thick,
        >=stealth,
        node distance=10mm and 12mm,
        layer/.style={
                draw,
                rounded corners=2pt,
                align=center,
                minimum height=14mm,
                minimum width=44mm,
                inner sep=4pt,
                font=\small,
            },
        example/.style={
                draw=none,
                align=center,
                font=\scriptsize\itshape,
            },
    ]
    \node[layer]                                    (count)   {Счётные\\(BoW, TF-IDF)};
    \node[layer, right=of count]                    (static)  {Статичные\\(word2vec, GloVe, fastText)};
    \node[layer, right=of static]                   (context) {Контекстные\\(BERT, XLM-R)};
    \node[layer, right=of context]                  (multi)   {Multilingual\\(SBERT-multi, LaBSE)};

    \node[example, below=4mm of count]   (ex1) {интерпретируемость,\\один язык};
    \node[example, below=4mm of static]  (ex2) {семантика слова,\\subword $\to$ OOV};
    \node[example, below=4mm of context] (ex3) {контекст,\\fine-tuning};
    \node[example, below=4mm of multi]   (ex4) {общее пространство\\ru + en};

    \draw[->] (count)   -- (static);
    \draw[->] (static)  -- (context);
    \draw[->] (context) -- (multi);
\end{tikzpicture}
}
    \caption{Типология векторных представлений текста, рассматриваемых в~настоящей
        работе: счётные, статичные дистрибутивные, контекстные и~многоязычные; в~каждой
        группе указаны характерные модели и~ключевые свойства}\label{fig:embedding-taxonomy}
\end{figure}

\section{Методы классификации текстов}\label{sec:ch1/clf}

Классификатор \( f_\theta \) в~постановке~\eqref{eq:decision-rule} обязан возвращать
оценки правдоподобия \( p_k(x) \) для каждого пользовательского класса. Этот раздел
рассматривает семейства классификаторов, востребованные в~прикладных работах по~почте,
с~акцентом на~три практически значимых критерия: совместимость с~ограниченной размеченной
выборкой, поддержку инкрементного обучения как~ответ на~дрейф концепций и~калибровку
вероятностей, без~которой правило с~порогом~\(\tau\) теряет смысл.

\subsection{Базовые семейства}\label{subsec:ch1/clf/baselines}

В~качестве контрольных моделей рассматриваются классические семейства, не~выбранные
как~целевое решение. Вероятностные и~линейные модели поверх TF-IDF "--- мультиномиальный
наивный Байес~\autocite{McCallum1998NaiveBayes}, линейные SVM и~логистическая
регрессия~\autocite{Joachims1998SVMTextCat,Fan2008LIBLINEAR,Pedregosa2011sklearn} "---
дёшевы и~устойчивы при~ограниченной разметке, но~не~дают калиброванной вероятности.
Онлайн-классификаторы "--- Passive-Aggressive~\autocite{Crammer2006PA}
и~Confidence-Weighted~\autocite{Dredze2008CW}, применённые
в~TaskTracer~\autocite{Keiser2009TaskTracer} "--- инкрементно обновляются при~дрейфе,
но~их уверенность относительна (margin), а~не~вероятность. Ансамбли деревьев "---
Random Forest~\autocite{Breiman2001RF}, XGBoost~\autocite{Chen2016XGBoost},
LightGBM~\autocite{Ke2017LightGBM} "--- сильны на~плотных признаках невысокой
размерности, но~проигрывают на~разреженных TF-IDF-векторах. Эти семейства служат
ориентиром сравнения, тогда как~целевая архитектура строится поверх многоязычных
эмбеддингов (раздел~\ref{subsec:ch1/repr/context}).

\subsection{Нейронный классификатор поверх эмбеддингов}\label{subsec:ch1/clf/neural}

Целевая архитектура\footnote{Компонент классификатора находится в~активной доработке;
ниже приводится его математическая постановка в~текущем, обобщённом виде.} "--- линейный
классификатор поверх векторного представления \( h = \varphi(x) \in \mathbb{R}^{d} \).
Логиты и~апостериорные оценки классов вычисляются как
\begin{align}
    z(x)   & = W h(x) + b, \qquad W \in \mathbb{R}^{K \times d},\ b \in \mathbb{R}^{K}, \label{eq:logits}                          \\
    p_k(x) & = \frac{\exp\bigl(z_k(x)\bigr)}{\sum_{j=1}^{K} \exp\bigl(z_j(x)\bigr)}, \qquad k = 1, \dots, K, \label{eq:softmax}
\end{align}
где \( h(x) \) "--- представление сообщения~\( x \) размерности~\( d \); \( W \) и~\( b \) "---
матрица весов и~вектор смещений линейного слоя; \( K \) "--- число пользовательских классов;
\( z_k(x) \) "--- \( k \)-й логит; \( p_k(x) \in (0, 1) \) "--- апостериорная оценка
класса~\( k \), причём \( \sum_{k} p_k(x) = 1 \). Параметры \( \theta = \{W, b, \dots\} \)
настраиваются минимизацией эмпирического кросс-энтропийного риска
\begin{equation}
    \label{eq:cross-entropy}
    \mathcal{L}(\theta) = -\frac{1}{N} \sum_{i=1}^{N} \sum_{k=1}^{K}
    \mathbb{1}\bigl[\,y_i = k\,\bigr] \log p_k(x_i),
\end{equation}
где \( N \) "--- размер обучающей выборки; \( (x_i, y_i) \) "--- размеченные пары
«сообщение"--~класс»; \( \mathbb{1}[\cdot] \) "--- индикатор. Оценки \( p_k(x) \) подаются
в~решающее правило с~порогом~\( \tau \) \eqref{eq:decision-rule}: при~\( \max_k p_k(x) < \tau \)
сообщение получает исход \texttt{other}. Калибровка вероятностей (temperature/Platt
scaling~\autocite{Guo2017Calibration}) и~softmax-эвристика выделения
out-of-distribution~\autocite{Hendrycks2017Baseline}, согласующие уровень уверенности
с~порогом и~отделяющие «незнакомое» от~«неоднозначного», относятся к~направлению
дальнейшей доработки компонента.

\section{Методы кластеризации для начального формирования классов}\label{sec:ch1/cluster}

В~постановке~\eqref{eq:label-set} множество пользовательских классов \(Y_{\text{user}}\)
не~задано заранее. Первичный набор классов формируется по~исторической выборке писем
кластеризацией в~признаковом пространстве \(\mathbb{R}^{d}\), задаваемом векторизатором
из~раздела~\ref{sec:ch1/repr}. Этот раздел рассматривает семейства методов кластеризации
с~точки зрения двух требований: способности работать без~априорного числа классов и~явного
выделения сообщений, не~попадающих ни~в~один из~устойчивых кластеров "--- естественного
прообраза категории \texttt{other}.

\subsection{Партиционные методы}\label{subsec:ch1/cluster/partition}

Базовая партиционная схема "--- алгоритм \(k\)-means~\autocite{MacQueen1967KMeans}
с~инициализацией \(k\)-means++~\autocite{Arthur2007KMeansPP}, минимизирующий суммарную
внутрикластерную дисперсию
\begin{equation}
    \label{eq:kmeans}
    \min_{\{C_1, \dots, C_k\}} \sum_{j=1}^{k} \sum_{x \in C_j}
    \bigl\lVert \varphi(x) - \mu_j \bigr\rVert^2,
    \qquad \mu_j = \frac{1}{\lvert C_j \rvert} \sum_{x \in C_j} \varphi(x),
\end{equation}
где \( C_j \) "--- \( j \)-й кластер (подмножество сообщений); \( \mu_j \) "--- его
центроид; \( k \) "--- фиксированное число кластеров; \( \lvert C_j \rvert \) "--- размер
кластера; \( \varphi(x) \) "--- векторное представление сообщения~\( x \);
\( \lVert \cdot \rVert \) "--- евклидова норма. Метод прост и~быстр, но~требует
фиксированного~\( k \), приписывает каждую точку к~ближайшему центру (без~выделения шума)
и~предполагает сферические кластеры. В~постановке с~неизвестным числом проектных категорий
и~разнородной массой шумовых сообщений эти ограничения существенны.

\subsection{Плотностные методы}\label{subsec:ch1/cluster/density}

Плотностные алгоритмы трактуют кластер как~область высокой плотности точек, разделённую
областями низкой плотности. DBSCAN~\autocite{Ester1996DBSCAN} задаёт ядровую точку
условием \( \lvert N_\varepsilon(x) \rvert \geqslant \mathit{minPts} \), где
\( N_\varepsilon(x) \) "--- множество точек в~радиусе~\( \varepsilon \) от~\( x \),
\( \mathit{minPts} \) "--- минимальное число соседей; не~прошедшие порог точки помечаются
как~шум. Иерархическое обобщение HDBSCAN~\autocite{Campello2013HDBSCAN,McInnes2017HDBSCAN}
устраняет зависимость от~единого~\( \varepsilon \), вводя ядровое расстояние
\( d_{\mathrm{core}}(x) \) "--- расстояние от~\( x \) до~его \( m \)-го ближайшего соседа
(\( m = \mathit{min\_samples} \)) "--- и~взаимную достижимость
\begin{equation}
    \label{eq:mreach}
    d_{\mathrm{mreach}}(a, b) = \max\bigl\{ d_{\mathrm{core}}(a),\,
    d_{\mathrm{core}}(b),\, d(a, b) \bigr\},
\end{equation}
где \( d(a, b) \) "--- базовое (евклидово) расстояние между точками~\( a \) и~\( b \);
\( d_{\mathrm{core}}(\cdot) \) "--- ядровое расстояние. В~разреженных областях ядровое
расстояние велико, поэтому одиночные точки «отодвигаются» друг от~друга
(рисунок~\ref{fig:hdbscan-mreach}). Кластеры строятся как~устойчивые поддеревья
минимального остова в~метрике \( d_{\mathrm{mreach}} \), а~устойчивость (persistence)
кластера~\( C \) измеряется как
\begin{equation}
    \label{eq:stability}
    S(C) = \sum_{x \in C} \bigl( \lambda_{\max}(x, C) - \lambda_{\mathrm{birth}}(C) \bigr),
    \qquad \lambda = \frac{1}{d_{\mathrm{mreach}}},
\end{equation}
где \( \lambda \) "--- уровень плотности (обратное расстояние);
\( \lambda_{\mathrm{birth}}(C) \) "--- уровень появления кластера~\( C \);
\( \lambda_{\max}(x, C) \) "--- уровень, на~котором точка~\( x \) покидает~\( C \);
\( S(C) \) "--- суммарная устойчивость (чем больше, тем стабильнее кластер). Отдельный
класс «шум» структурно сопоставим с~категорией \texttt{other} из~правила~\eqref{eq:decision-rule}:
сообщения, не~попавшие ни~в~один устойчивый кластер на~этапе начального формирования
классов, естественно отнести в~буфер обратной связи, а~не~принудительно
распределить по~пользовательским классам (рисунок~\ref{fig:clustering-outliers}).

\begin{figure}[h]
    \centering
    \includegraphics[width=0.6\textwidth]{hdbscan_mreach}
    \caption{Взаимная достижимость HDBSCAN на~синтетических точках: вокруг \( a \)
        и~\( b \) показаны окружности ядрового расстояния \( d_{\mathrm{core}} \)
        (радиус "--- до~\( m \)-го соседа); для точки~\( b \) в~разреженной области
        ядровое расстояние велико, поэтому \( d_{\mathrm{mreach}}(a, b) \) определяется
        им, а~не~прямым расстоянием \( d(a, b) \)}\label{fig:hdbscan-mreach}
\end{figure}

\subsection{Снижение размерности}\label{subsec:ch1/cluster/dimred}

Контекстные эмбеддинги размещают сообщения в~пространстве сотен измерений, где
плотностные алгоритмы теряют устойчивость. Поэтому перед HDBSCAN применяется нелинейное
снижение размерности UMAP~\autocite{McInnes2018UMAP}: метод строит взвешенный граф
\( k \)-ближайших соседей в~исходном пространстве и~ищет низкоразмерное вложение,
сохраняющее его локальную структуру. В~настоящей работе UMAP включён как~подготовительный
шаг кластеризации, после которого HDBSCAN работает в~пространстве небольшой размерности.

\subsection{Метрики качества кластеризации}\label{subsec:ch1/cluster/metrics}

Основной внутренней мерой качества в~работе служит плотностной индекс валидности
DBCV~\autocite{Moulavi2014DBCV}, согласованный с~метрикой взаимной достижимости HDBSCAN.
Для кластера~\( C_i \) вводятся внутрикластерная разреженность \( \mathrm{DSC}(C_i) \) "---
наибольшее ребро минимального остова \( C_i \) в~метрике \( d_{\mathrm{mreach}} \) "---
и~межкластерная разделимость \( \mathrm{DSPC}(C_i, C_j) \) "--- наименьшая взаимная
достижимость между точками \( C_i \) и~\( C_j \); валидность кластера
\begin{equation}
    \label{eq:dbcv-cluster}
    V(C_i) = \frac{\min_{j \neq i} \mathrm{DSPC}(C_i, C_j) - \mathrm{DSC}(C_i)}
    {\max\bigl\{ \min_{j \neq i} \mathrm{DSPC}(C_i, C_j),\ \mathrm{DSC}(C_i) \bigr\}},
\end{equation}
а~итоговый индекс "--- средневзвешенное по~размерам кластеров
\begin{equation}
    \label{eq:dbcv}
    \mathrm{DBCV} = \sum_{i} \frac{\lvert C_i \rvert}{N}\, V(C_i) \in [-1, 1],
\end{equation}
где \( \mathrm{DSC}(C_i) \) "--- плотностная разреженность внутри кластера~\( C_i \);
\( \mathrm{DSPC}(C_i, C_j) \) "--- плотностная разделимость кластеров~\( C_i \) и~\( C_j \);
\( V(C_i) \in [-1, 1] \) "--- валидность кластера~\( C_i \); \( \lvert C_i \rvert \) "---
его размер; \( N \) "--- общее число точек. Большие значения \( \mathrm{DBCV} \) отвечают
плотным и~хорошо разделённым кластерам (рисунок~\ref{fig:dbcv-validity}). При~автоматическом
подборе параметров HDBSCAN (Optuna~\autocite{Akiba2019Optuna}) целевая функция объединяет
DBCV с~устойчивостью~\eqref{eq:stability} и~средней уверенностью принадлежности при~штрафе
за~долю шума, а~поиск ведётся в~области допустимых решений (число кластеров, доля шума,
медианный размер кластера). Классические индексы "--- коэффициент
силуэта~\autocite{Rousseeuw1987Silhouette} и~индекс Дэвиса"--~Боулдина~\autocite{Davies1979DBI} "---
предполагают выпуклые кластеры и~используются лишь как~вспомогательный ориентир;
окончательное решение о~принятии кластеров требует ручной проверки кандидатов.

\begin{figure}[h]
    \centering
    \includegraphics[width=0.6\textwidth]{dbcv_density_validity}
    \caption{Плотностная валидность DBCV на~синтетических данных: рёбра показывают
        минимальные остовы двух кластеров в~метрике взаимной достижимости; выделены
        внутрикластерная разреженность \( \mathrm{DSC} \) (наибольшее внутреннее ребро)
        и~межкластерная разделимость \( \mathrm{DSPC} \) (наименьшее ребро между
        кластерами), из~которых складывается индекс}\label{fig:dbcv-validity}
\end{figure}

\begin{figure}[h]
    \centering
    \begin{tikzpicture}[
        thick,
        >=stealth,
        every node/.style={font=\scriptsize},
        cluster1/.style={circle, draw=black, fill=black!10, inner sep=1pt, minimum size=2.6mm},
        cluster2/.style={circle, draw=black, fill=black!30, inner sep=1pt, minimum size=2.6mm},
        cluster3/.style={circle, draw=black, fill=black!60, text=white, inner sep=1pt, minimum size=2.6mm},
        noise/.style={circle, draw=black, fill=white, inner sep=1pt, minimum size=2.4mm, dashed},
        panel/.style={draw, rounded corners=2pt, minimum width=56mm, minimum height=42mm},
    ]

    \begin{scope}[local bounding box=left]
        \node[panel] (Lframe) at (0,0) {};
        \node[anchor=south, font=\scriptsize, yshift=1.2mm] at (Lframe.north) {а) k-means: все точки приписаны};
        % три кластера + «выбросы», но все покрашены
        \foreach \x/\y in {-1.6/0.7, -1.2/1.1, -1.8/0.4, -1.4/0.2, -1.1/0.9}
            \node[cluster1] at (\x,\y) {};
        \foreach \x/\y in {0.6/0.5, 1.0/0.9, 0.4/1.0, 0.9/0.3, 1.2/0.6}
            \node[cluster2] at (\x,\y) {};
        \foreach \x/\y in {-0.4/-1.0, -0.1/-1.4, 0.2/-1.1, -0.3/-0.7}
            \node[cluster3] at (\x,\y) {};
        % «истинные выбросы», тоже покрашенные (искусственно отнесённые)
        \foreach \x/\y/\c in {1.7/-1.3/cluster3, -1.8/-1.2/cluster1, 1.8/1.4/cluster2}
            \node[\c] at (\x,\y) {};
    \end{scope}

    \begin{scope}[xshift=72mm, local bounding box=right]
        \node[panel] (Rframe) at (0,0) {};
        \node[anchor=south, font=\scriptsize, yshift=1.2mm] at (Rframe.north) {б) HDBSCAN: часть точек "--- шум};
        % те же три кластера
        \foreach \x/\y in {-1.6/0.7, -1.2/1.1, -1.8/0.4, -1.4/0.2, -1.1/0.9}
            \node[cluster1] at (\x,\y) {};
        \foreach \x/\y in {0.6/0.5, 1.0/0.9, 0.4/1.0, 0.9/0.3, 1.2/0.6}
            \node[cluster2] at (\x,\y) {};
        \foreach \x/\y in {-0.4/-1.0, -0.1/-1.4, 0.2/-1.1, -0.3/-0.7}
            \node[cluster3] at (\x,\y) {};
        % выбросы — отдельный стиль (шум)
        \foreach \x/\y in {1.7/-1.3, -1.8/-1.2, 1.8/1.4}
            \node[noise] at (\x,\y) {};
    \end{scope}
\end{tikzpicture}

    \caption{Сравнение партиционной и~плотностной кластеризации: \emph{а)}~\(k\)-means
        приписывает каждую точку к~ближайшему центру, не~различая шум; \emph{б)}~HDBSCAN
        выделяет три плотных кластера и~помечает оставшиеся точки как~шум (пунктирные
        обводки), что естественно сопоставимо с~категорией~\texttt{other}}\label{fig:clustering-outliers}
\end{figure}

Таким образом, по~требованиям раздела предпочтительны плотностные методы: HDBSCAN
выделяет устойчивые кластеры переменной плотности и~помечает остальные точки как~шум,
сопоставимый с~отказом \texttt{other}, а~партиционный \( k \)-means сохраняется
как~базовая линия.

\section{Сравнительный анализ существующих систем}\label{sec:ch1/systems}

Свойства задачи из~раздела~\ref{sec:ch1/domain} "--- персонализированные классы,
многоязычность, локальность исполнения и~учёт дрейфа концепций "--- вместе с~требованием
категории \texttt{other} и~обратной связи~\eqref{eq:decision-rule} образуют оси
сравнения, по~которым ниже сопоставляются четыре характерные системы организации
входящей корреспонденции.

Focused~Inbox (Outlook)~\autocite{Microsoft2024FocusedInbox} и~Priority~Inbox
(Gmail)~\autocite{Aberdeen2010PriorityInbox} персонализируют важность письма, но~сводят
результат к~бинарной шкале «важное/прочее», исполняются на~серверах сервиса
и~не~допускают произвольных тематических классов; их адаптация к~пользователю
не~декларируется как~явный механизм учёта дрейфа концепций, а~многоязычность
не~специфицирована. TaskTracer~\autocite{Keiser2009TaskTracer} предсказывает проектную
принадлежность письма онлайн-классификаторами (наивный Байес, Passive-Aggressive,
Confidence-Weighted) и~приближается к~многоклассовой постановке~\eqref{eq:classifier},
но~работает только с~англоязычной перепиской и~без~категории \texttt{other}.
EmFORE~\autocite{Singh2023EmFORE} индуктивно выводит символьные правила распределения
по~папкам с~явной адаптацией к~дрейфу признаков, опираясь на~структурные признаки
(отправитель, шаблоны темы и~заголовков), но~без~кластеризации для формирования классов
и~без~активного буфера \texttt{other}; многоязычное применение не~декларируется.

Сопоставление перечисленных систем по~осям, существенным для постановки настоящей
работы, приведено в~таблице~\ref{tab:systems-compare}. Ручные правила почтового
клиента и~пользовательские фильтры Sieve рассматриваются в~разделе~\ref{sec:ch1/domain}
и~в~таблице отдельной строкой не~представлены.

\begin{table}[h]
    \centering
    \begin{threeparttable}
        \caption{Сравнение существующих систем организации входящей почты с~целевыми
            требованиями настоящей работы}\label{tab:systems-compare}
        \small\setlength{\tabcolsep}{4pt}
        \begin{tabular}{| p{26mm} | c | c | c | c | c | >{\raggedright\arraybackslash}p{30mm} |}
            \hline
            Система                      & \makecell{Перс.\\классы\tnote{1}} & \makecell{Много-\\язычн.} & \makecell{Лок.\\испол.} & \makecell{\texttt{other}\tnote{2}} & \makecell{Дрейф\\конц.} & \makecell[l]{Источник\\классов} \\
            \hline
            Focused Inbox~\autocite{Microsoft2024FocusedInbox}
                                         & \(-\)\tnote{3}               & н/д          & \(-\)        & \(-\)                   & \(\pm\)\tnote{4} & встроенная бинарная \\
            \hline
            Priority Inbox~\autocite{Aberdeen2010PriorityInbox}
                                         & \(-\)\tnote{3}               & н/д          & \(-\)        & \(-\)                   & \(\pm\)\tnote{4} & встроенная бинарная \\
            \hline
            TaskTracer~\autocite{Keiser2009TaskTracer}
                                         & \(+\)                        & \(-\)        & \(+\)        & \(-\)                   & \(+\)         & пользова\-тельские проекты \\
            \hline
            EmFORE~\autocite{Singh2023EmFORE}
                                         & \(+\)                        & \(-\)        & \(\pm\)      & \(-\)                   & \(+\)         & папки и~индуктивные правила \\
            \hline
            \textbf{Настоящая работа}    & \(+\)                        & \(+\)        & \(+\)        & \(+\)                   & \(+\)         & кластеризация архива \\
            \hline
        \end{tabular}
        \begin{tablenotes}\footnotesize
            \item[1] Произвольный набор тематических классов, заданный пользователем.
            \item[2] Явная категория для сообщений вне устойчивых классов с~накоплением для переобучения.
            \item[3] Поддерживается только бинарное разделение «важное/прочее».
            \item[4] Адаптация по~неявным сигналам пользователя; явный механизм учёта дрейфа концепций отсутствует.
        \end{tablenotes}
    \end{threeparttable}
\end{table}

Из~таблицы~\ref{tab:systems-compare} следует, что ни~одна из~рассмотренных систем
не~закрывает все оси одновременно: системы крупных почтовых сервисов поддерживают
дрейф концепций, но~не~дают произвольных классов и~требуют передачи данных вовне; работы
TaskTracer и~EmFORE дают персонализированные классы, но~не~рассчитаны на~многоязычность
и~явное обращение с~категорией \texttt{other}. Целевая строка таблицы фиксирует
требования, которым должна удовлетворять система настоящей работы.

\section{Постановка задачи разработки}\label{sec:ch1/dev}

\subsection{Резюме предметной области}\label{subsec:ch1/dev/summary}

Постановка~\eqref{eq:classifier}--\eqref{eq:decision-rule} вместе со~свойствами
из~раздела~\ref{sec:ch1/domain} и~метриками раздела~\ref{subsec:ch1/task/metrics}
формирует следующую совокупность требований к~системе: персонализированные
тематические классы, формируемые по~исторической выборке; единое признаковое
пространство для русско-английской переписки; локальное исполнение всех компонентов;
учёт дрейфа концепций с~явной категорией~\texttt{other} как~буфером обратной связи.
Сопоставление с~существующими системами из~таблицы~\ref{tab:systems-compare}
показывает, что ни~одно из~готовых решений не~удовлетворяет всем требованиям
одновременно, что и~определяет содержание дальнейшей работы.

\subsection{Список задач разработки}\label{subsec:ch1/dev/tasks}

В~рамках работы решаются следующие задачи разработки.

\begin{enumerate}
    \item Нормализация входящих сообщений: построение конвейера, переводящего MIME-сообщение
        в~канонический контракт \texttt{normalized\_text}, включающий поля заголовка,
        очищенный основной текст и~идентификатор языка.
    \item Векторизация: реализация локального многоязычного векторизатора~\(\varphi\)
        для русско-английской переписки, не~требующего передачи текста писем во~внешние
        сервисы; конкретный выбор семейства методов "--- глава~\ref{ch:ch2}.
    \item Начальное формирование пользовательских классов: процедура кластеризации архива писем
        в~пространстве \(\varphi\) с~явным выделением шумовых сообщений и~подготовкой
        кандидатов для ручной проверки.
    \item Обучение классификатора: реализация классификатора \(f_\theta\), отдающего
        калиброванные оценки \(p_k(x)\), и~интеграция решающего правила~\eqref{eq:decision-rule}
        с~порогами по~уверенности и~отнесением низкоуверенных сообщений в~\texttt{other}.
    \item Цикл обратной связи: накопление сообщений категории \texttt{other},
        периодическая ручная проверка и~пороговый триггер переобучения как~механизм
        реакции на~дрейф концепций~\autocite{Gama2014ConceptDrift}.
    \item Инкрементное обновление и~версионирование: хранение версий модели вместе
        с~метриками macro-F1, взвешенного F1 и~покрытия, переключение на~новую версию
        по~критериям активации (quality gates).
    \item Эмпирическая оценка: проведение экспериментов на~архиве писем,
        репрезентативном для целевого пользователя, с~отчётом по~метрикам
        раздела~\ref{subsec:ch1/task/metrics}.
\end{enumerate}

\subsection{Критерии успеха}\label{subsec:ch1/dev/success}

Система считается решающей поставленную задачу при~одновременном выполнении следующих
критериев.

\begin{enumerate}
    \item Качество классификации на~тестовом срезе достигает значения
        \( \mathrm{macro\text{-}F_1} \geqslant \fixme{\tau_{\mathrm{F1}}} \) при~покрытии
        не~ниже \( \fixme{c_{0}} \) от~общего объёма входящих писем; конкретные
        числовые пороги фиксируются по~результатам экспериментов в~главе~\ref{ch:ch3}.
    \item Среднее время инференса одного сообщения на~целевой конфигурации
        не~превышает \( \fixme{\tau_{\mathrm{lat}}} \)~мс на~CPU.
    \item Архитектура и~процедуры эксплуатации формально гарантируют отсутствие
        передачи текста писем во~внешние сервисы.
    \item Реализован полный цикл «\texttt{other}~$\to$~ручная проверка~$\to$~переобучение
        $\to$~активация», подтверждённый сценариями использования в~главе~\ref{ch:ch3}.
\end{enumerate}

Архитектурная реализация перечисленных задач разработки, выбор конкретных моделей
и~алгоритмов, а~также организация цикла обратной связи рассматриваются
в~главе~\ref{ch:ch2}.


\FloatBarrier
